% !PRINTED 99
\weberpage{99}

From the equation

\weq{8}{\alpha\delta - \beta\gamma = 1}

it follows that neither $\alpha, \beta$ nor $\alpha, \gamma$, nor $\delta, \beta$, nor $\delta, \gamma$ can have a common factor. But if $\alpha, \beta$ have been assumed arbitrarily without a common divisor, then $\gamma, \delta$ can still be determined from (8) in infinitely many ways. If $\gamma, \delta$ is one of these determinations, then they are all contained in the form

\weq{9}{\gamma + \lambda\alpha, \quad \delta + \lambda\beta}

in which $\lambda$ is an arbitrary integer (Bd. I, \S\ 126).

\begin{center}
\textbf{\S\ 29.\quad\gesperrt{Composition of the transformations from simpler ones.}}
\end{center}

In the system $\mathfrak{S}$ of all transformations $S$ there is contained a system $\mathfrak{S}_0$ which consists of all those transformations $S_0$ whose second transformation number is $b = 0$, while $a$ and $\partial$ are \gesperrt{positive}:

\weq{1}{S_0 = \begin{pmatrix} a, & 0 \\ c, & \partial \end{pmatrix}.}

When two $S_0$ are composed there arises again an $S_0$; but $\mathfrak{S}_0$ is nevertheless just as little a group as $\mathfrak{S}$.

One can reduce any arbitrary transformation
\[
S = \begin{pmatrix} p, & q \\ r, & s \end{pmatrix}
\]
to an $S_0$ by a composition $LS$. For if

\weq{2}{\begin{pmatrix} \alpha, & \beta \\ \gamma, & \delta \end{pmatrix}\begin{pmatrix} p, & q \\ r, & s \end{pmatrix} = \begin{pmatrix} a, & 0 \\ c, & \partial \end{pmatrix}}

is to hold, then $\alpha, \beta$ must satisfy the condition:
\[
\alpha q + \beta s = 0,
\]
and accordingly, if $\partial$ is the greatest common divisor of $q$ and $s$, one sets
\[
\partial\alpha = s, \quad \partial\beta = -q,
\]
and, after $\alpha$ and $\beta$ have hereby been ascertained as relatively prime numbers, one determines $\gamma$ and $\delta$ from formula (8), \S\ 28. Then (2) is satisfied if
\[
a\partial = n, \quad c = \gamma p + \delta r
\]
is put; $a$ and $\partial$ are hereby uniquely determined, but $c$ can, for another choice of $\gamma$ and $\delta$, be replaced by $c+\lambda a$. One can therefore dispose of $\lambda$ so that $c$ lies in the sequence
%CONT>>
